\cleardoublepage
\chapter{KODE PROGRAM}
\label{chap:lampiran-a}
Pada umumnya, kode program tidak perlu disertakan di dalam laporan tugas akhir karena kode program biasanya sangat panjang dan sudah disimpan dalam bentuk berkas-berkas elektronik terpisah di repositori daring. Namun, bagian kode program singkat yang penting untuk memahami implementasi sistem disertakan berikut ini.

\section{Kode Program \textit{Frontend}}

\begin{lstlisting}[style=javascriptstyle, caption={Implementasi Alur \textit{Login} dan Penyimpanan Sesi}, label={lst:login}]
const handleSubmit = async (e) => {
  e.preventDefault();
  setLoading(true);
  try {
    const data = await loginUser(form.email, form.password);
    saveSession(data);
    await loadProgressFromDb(data.token);
    router.push('/dashboard');
  } catch (err) {
    setError(err.message);
  } finally {
    setLoading(false);
  }
};

// auth.js
export function saveSession(data) {
  localStorage.setItem('token', data.token);
  localStorage.setItem('user', JSON.stringify(data.user));
}
\end{lstlisting}

\begin{lstlisting}[style=javascriptstyle, caption={Fungsi \texttt{buildInsights} untuk Rekomendasi Topik}, label={lst:beranda-insights}]
function buildInsights(rows) {
  const list = [];
  rows.forEach(({ topic, d, secs }) => {
    if (secs === 0) return;
    if (secs === 3) {
      const missing = SECTIONS.find((s) => !d[s]);
      list.push({ priority: 1, topic,
        message: `Hampir selesai di ${topic.shortTitle} — tinggal ${SECTION_LABELS[missing]}.`,
        cta: 'Selesaikan' });
    } else if (d.materi && d.contoh && !d.latihan) {
      list.push({ priority: 2, topic,
        message: `Sudah baca materi ${topic.shortTitle}, tapi belum mencoba latihan.`,
        cta: 'Coba latihan' });
    } else if (d.materi && !d.contoh) {
      list.push({ priority: 3, topic,
        message: `Pelajari contoh soal ${topic.shortTitle} untuk memperkuat pemahaman.`,
        cta: 'Lihat contoh' });
    }
  });
  return list.sort((a, b) => a.priority - b.priority).slice(0, 3);
}
\end{lstlisting}

\begin{lstlisting}[style=javascriptstyle, caption={Perhitungan Progres Topik dari \texttt{localStorage}}, label={lst:topik-progress}]
function calcProgress(key) {
  try {
    const d = JSON.parse(localStorage.getItem(key)) ?? {};
    return ['materi', 'contoh', 'latihan', 'ringkasan']
      .filter((s) => d[s]).length * 25;
  } catch { return 0; }
}

// Inisialisasi saat komponen mount
useEffect(() => {
  setPengantarProgress(calcProgress('asd_progress_pengantar'));
  setAdtSederhanaProgress(calcProgress('asd_progress_adt_sederhana'));
  // ... topik lainnya
}, []);
\end{lstlisting}

\begin{lstlisting}[style=javascriptstyle, caption={Fungsi Penyimpanan Progres Sesi ke \texttt{localStorage} dan \textit{Backend}}, label={lst:materi-progress}]
function saveProgress(updates) {
  const data = { ...readProgress(), ...updates };
  localStorage.setItem(PROGRESS_KEY, JSON.stringify(data));
  syncTopicProgress(TOPIC_SLUG, data);
}

const handleComplete = (tab) => {
  const key = TAB_KEYS[tab];          // 'materi' | 'contoh' | 'latihan' | 'ringkasan'
  if (!key || completed[key]) return;
  saveProgress({ [key]: true });
  completedRef.current = { ...completedRef.current, [key]: true };
  forceUpdate((n) => n + 1);
};
\end{lstlisting}

\begin{lstlisting}[style=javascriptstyle, caption={Mekanisme \textit{Toggle} Hint dan Pembahasan pada Halaman Contoh}, label={lst:contoh-toggle}]
const [showHint, setShowHint]       = useState(false);
const [showJawaban, setShowJawaban] = useState(false);

// Tombol Hint
<button onClick={() => setShowHint(!showHint)}>
  {showHint ? "v" : ">"} Hint
</button>
{showHint && <div>...petunjuk pendekatan...</div>}

// Tombol Pembahasan
<button onClick={() => setShowJawaban(!showJawaban)}>
  {showJawaban ? "v" : ">"} Lihat Pembahasan
</button>
{showJawaban && <div>...algoritma, notasi, kode C, trace...</div>}
\end{lstlisting}

\begin{lstlisting}[style=javascriptstyle, caption={Alur \textit{Generate} Soal Adaptif dan Evaluasi per Soal}, label={lst:latihan-ai}]
const generateSoal = useCallback(async (kelemahan = []) => {
  setFase("loading");
  const res = await fetch("/api/latihan/generate", {
    method: "POST",
    body: JSON.stringify({ jumlah: 5, kelemahan }),
  });
  const data = await res.json();
  setSoalList(data.soal);
  setFase("latihan");
}, []);

const evaluasiSoal = async () => {
  const res = await fetch("/api/latihan/evaluasi", {
    method: "POST",
    body: JSON.stringify({ soal, jawaban: jawaban[soal.id] }),
  });
  const data = await res.json();
  setFeedbackMap((prev) => ({ ...prev, [soal.id]: data }));
};

// Saat generate baru, kirim konsep lemah dari sesi sebelumnya
const handleGenerateBaru = () => {
  const kelemahan = [
    ...new Set(Object.values(feedbackMap).flatMap((f) => f.konsep_lemah ?? [])),
  ];
  generateSoal(kelemahan);
};
\end{lstlisting}

\begin{lstlisting}[style=javascriptstyle, caption={Fungsi Perhitungan \textit{Streak} Belajar Harian}, label={lst:progress-streak}]
function updateStreak() {
  const today = new Date().toISOString().split('T')[0];
  const lastVisit = localStorage.getItem('asd_last_visit');
  let streak = parseInt(localStorage.getItem('asd_streak') ?? '0', 10);
  if (!lastVisit) {
    localStorage.setItem('asd_last_visit', today);
    localStorage.setItem('asd_streak', '1');
    return 1;
  }
  if (lastVisit === today) return streak || 1;
  const diff = Math.round(
    (new Date(today) - new Date(lastVisit)) / 86400000
  );
  streak = diff === 1 ? streak + 1 : 1;
  localStorage.setItem('asd_last_visit', today);
  localStorage.setItem('asd_streak', String(streak));
  return streak;
}
\end{lstlisting}

\begin{lstlisting}[style=javascriptstyle, caption={Fungsi \texttt{refreshSession} dan Alur Penyimpanan Perubahan Profil}, label={lst:pengaturan-profile}]
const refreshSession = (updatedUser) => {
  const s = getSession();
  const newUser = { ...s.user, ...updatedUser };
  saveSession({ token: s.token, user: newUser });   // perbarui localStorage
  setSession((prev) => ({ ...prev, user: newUser }));
};

const saveName = async () => {
  if (!nameDraft.trim())
    return setStatus({ field: 'name', error: 'Nama tidak boleh kosong.' });
  setLoading(true);
  try {
    const updated = await updateProfile(session.token, { name: nameDraft.trim() });
    refreshSession(updated);
    setEditing(null);
  } catch (err) {
    setStatus({ field: 'name', error: err.message });
  } finally { setLoading(false); }
};
\end{lstlisting}

\section{Kode Program \textit{Backend}}

\begin{lstlisting}[style=javascriptstyle, caption={Implementasi \textit{Password Hashing} dan \textit{Dependency} Autentikasi}, label={lst:backend-auth}]
# security.py
def hash_password(password: str) -> str:
    salt = secrets.token_hex(PBKDF2_SALT_BYTES)   # 16 byte acak
    h = hashlib.pbkdf2_hmac("sha256", password.encode(),
                             salt.encode(), 100_000).hex()
    return f"pbkdf2_sha256$100000${salt}${h}"

# routes.py
def get_current_user(
    authorization: str = Header(...),
    db: Session = Depends(get_db),
) -> User:
    token = authorization.removeprefix("Bearer ")
    payload = decode_access_token(token, settings.secret_key)
    user_id = int(payload["sub"])
    user = db.execute(
        select(User).where(User.user_id == user_id)
    ).scalar_one_or_none()
    if user is None:
        raise HTTPException(status_code=404, detail="Pengguna tidak ditemukan.")
    return user
\end{lstlisting}

\begin{lstlisting}[style=javascriptstyle, caption={Implementasi \textit{Upsert} Progres Belajar per Topik}, label={lst:backend-progress}]
@router.put("/progress/{topic_slug}", response_model=ProgressResponse)
async def upsert_progress(
    topic_slug: str,
    payload: ProgressUpdateRequest,
    current_user: User = Depends(get_current_user),
    db: Session = Depends(get_db),
) -> ProgressResponse:
    row = db.execute(
        select(Progress).where(
            Progress.user_id == current_user.user_id,
            Progress.topic_slug == topic_slug,
        )
    ).scalar_one_or_none()

    if row is None:                          # belum ada — buat baru
        row = Progress(user_id=current_user.user_id, topic_slug=topic_slug)
        db.add(row)

    row.materi    = payload.materi           # perbarui status 4 sesi
    row.contoh    = payload.contoh
    row.latihan   = payload.latihan
    row.ringkasan = payload.ringkasan
    row.last_accessed = datetime.now(timezone.utc)
    db.commit()
    db.refresh(row)
    return ProgressResponse.model_validate(row)
\end{lstlisting}

\begin{lstlisting}[style=javascriptstyle, caption={Implementasi \textit{Endpoint} Generate dan Evaluasi Soal Adaptif}, label={lst:backend-ai}]
@router.post("/latihan/{topic_slug}/generate")
async def generate_latihan(
    topic_slug: str, payload: GenerateSoalRequest,
) -> dict:
    result = await generate_soal(
        topic_slug=topic_slug,
        jumlah=payload.jumlah,
        kelemahan=payload.kelemahan,      # konsep lemah dari sesi sebelumnya
        tipe_paksa=payload.tipe_paksa,
        topik_referensi=payload.topik_referensi,
    )
    return result

@router.post("/latihan/{topic_slug}/evaluasi")
async def evaluasi_latihan(
    topic_slug: str, payload: EvaluasiSoalRequest,
) -> dict:
    # model mengembalikan: skor, nilai, metrik, komentar, konsep_lemah
    return await evaluasi_soal(
        topic_slug=topic_slug,
        soal=payload.soal,
        jawaban=payload.jawaban,
    )
\end{lstlisting}
